\chapter{Sheaf condition along the structure-sheaf forget composite}
\label{chap:Cohomology_SheafCompose}

The structure sheaf of a scheme is naturally a sheaf of commutative rings. Its underlying additive groups also form a sheaf, because the relevant forgetful functor preserves the limits that express the sheaf condition.

\section{Statement}

We record this preservation along the forgetful composite
\(\mathrm{CommRing} \to \mathrm{Ring} \to \mathrm{Ab}\).

\begin{theorem}[Sheaf condition transports along the structure-sheaf forget composite]\leanok
  \label{thm:HasSheafCompose_forget}
  \dcref{custom}
  \lean{AlgebraicGeometry.Cohomology.instHasSheafCompose_forget_CommRing_AddCommGrp}
  Let \(X\) be a topological space. Then the Grothendieck topology of opens of \(X\) transports sheaves along the forgetful composite from commutative rings to abelian groups.
\end{theorem}

\begin{proof}


\leanok
  Both forgetful functors in the composite \(\mathrm{CommRing} \to \mathrm{Ring} \to \mathrm{Ab}\) create (in particular preserve) all small limits. The sheaf condition for the topology of opens is expressed as preservation of certain multiequalizer limits. Composing two limit-preserving functors gives a limit-preserving functor, so post-composition with the forget composite preserves the sheaf condition.
\end{proof}
